% Part: incompleteness
% Chapter: representability-in-q
% Section: beta-function

\documentclass[../../../include/open-logic-section]{subfiles}

\begin{document}

\olfileid[id]{inc}{req}{bet}
\olsection{Lema Fungsi Beta}


Untuk menunjukkan bahwa kita dapat melaksanakan rekursi primitif jika
penjumlahan, perkalian, dan $\Char{=}$ tersedia, kita perlu mengembangkan
fungsi-fungsi yang menangani barisan. (Jika perpangkatan juga tersedia, tugas
kita akan lebih mudah.) Ketika rekursi primitif tersedia, kita dapat
mendefinisikan hal-hal seperti ``bilangan prima ke-$n$,'' lalu memilih
pengodean yang cukup langsung. Namun, di sini rekursi primitif tidak
tersedia---bahkan, kita ingin menunjukkan bahwa rekursi primitif dapat
dilakukan dengan menggunakan pencarian tak berbatas---sehingga kita harus lebih cerdik.

\begin{lem}
\ollabel{lem:beta}
Terdapat fungsi $\beta(d,i)$ sedemikian sehingga, untuk setiap barisan $a_0$,
\dots,~$a_n$, terdapat bilangan~$d$ sedemikian sehingga untuk setiap
$i \le n$, $\beta(d,i) = a_i$. Selain itu, $\beta$ dapat didefinisikan dari
fungsi-fungsi dasar hanya dengan komposisi dan pencarian tak berbatas pada
fungsi reguler.
\end{lem}

Bayangkan $d$ sebagai kode barisan $\tuple{a_0, \dots, a_n}$ dan
$\beta(d,i)$ sebagai fungsi yang mengembalikan elemen ke-$i$. (Perhatikan
bahwa ``pengodean'' ini \emph{tidak} menggunakan pengodean dengan pangkat
bilangan prima yang sudah kita kenal!) Lema tersebut cukup minimal; lema itu
tidak mengatakan bahwa kita dapat menggabungkan barisan, menambahkan elemen,
atau bahkan \emph{menghitung}~$d$ dari $a_0$, \dots,~$a_n$ dengan fungsi yang
dapat didefinisikan melalui komposisi dan pencarian tak berbatas pada fungsi
reguler. Lema itu hanya mengatakan bahwa terdapat suatu fungsi ``pengurai
kode'' sedemikian sehingga setiap barisan ``dikodekan.''

Penggunaan notasi $\beta$ berasal dari G\"odel. Sekali lagi, bagian sulit
dalam membuktikan lema tersebut ialah mendefinisikan~$\beta$ yang sesuai
dengan sumber daya yang tampaknya terbatas, yakni hanya menggunakan komposisi
dan pencarian tak berbatas---namun, kita diperbolehkan menggunakan penjumlahan,
perkalian, dan~$\Char{=}$. Ada berbagai cara untuk membuktikan lema ini,
tetapi salah satu yang paling jernih tetap metode asli G\"odel, yang
menggunakan fakta teori bilangan bernama Teorema Sunzi (secara tradisional,
``Teorema Sisa Tiongkok'').

\begin{defn}
Dua bilangan asli $a$ dan $b$ disebut \emph{relatif prima} jika dan hanya
jika faktor persekutuan terbesarnya ialah~$1$; dengan kata lain, keduanya tidak
memiliki pembagi bersama lain.
\end{defn}

\begin{defn}
Bilangan asli $a$ dan $b$ disebut \emph{kongruen modulo~$c$},
$a \equiv b \mod c$, jika dan hanya jika $c \mid (a-b)$, yakni $a$ dan $b$
memiliki sisa yang sama ketika dibagi dengan~$c$.
\end{defn}

Berikut Teorema Sunzi:
\begin{thm}
Andaikan $x_0$, \dots,~$x_n$ relatif prima sepasang-sepasang. Misalkan
$y_0$, \dots,~$y_n$ sebarang bilangan. Maka terdapat bilangan $z$ sedemikian
sehingga
\begin{align*}
z & \equiv y_0 \mod x_0 \\
z & \equiv y_1 \mod x_1 \\
& \vdots  \\
z & \equiv y_n \mod x_n.
\end{align*}
\end{thm}

Kita akan menggunakan Teorema Sunzi sebagai berikut: jika $x_0$,
\dots,~$x_n$ masing-masing lebih besar daripada $y_0$, \dots,~$y_n$, kita
dapat menggunakan $z$ untuk mengodekan barisan $\tuple{y_0, \dots,y_n}$.
Untuk memperoleh kembali~$y_i$, kita hanya perlu membagi $z$ dengan~$x_i$
dan mengambil sisanya. Untuk menggunakan pengodean ini, kita perlu menemukan
nilai-nilai yang sesuai bagi $x_0$, \dots,~$x_n$.

Beberapa pengamatan akan membantu kita dalam hal ini. Diberikan
$y_0$, \dots,~$y_n$, misalkan
\begin{align*}
j &= \max(n, y_0 + 1, \dots, y_n + 1), \\
m &= \lcm(1,\dots,j),
\end{align*}
dan misalkan
\begin{align*}
x_0 & = 1 + m \\
x_1 & = 1 + 2 \cdot m \\
x_2 & = 1 + 3 \cdot m \\
& \vdots  \\
x_n & = 1 + (n+1) \cdot m
\end{align*}
Maka dua hal berikut benar:
\begin{enumerate}
\item\ollabel{rel-prime} $x_0,\dots,x_n$ relatif prima sepasang-sepasang.
\item\ollabel{less} Untuk setiap $i$, $y_i < x_i$.
\end{enumerate}
Untuk melihat bahwa \olref{rel-prime} benar, perhatikan bahwa jika $p$ adalah
bilangan prima dan $p \mid x_i$ serta $p \mid x_k$, maka
$p \mid 1 + (i+1) m$ dan $p \mid 1 + (k+1) m$. Akan tetapi, $p$ kemudian
membagi selisih keduanya,
\[
(1 + (i+1)m) - (1+ (k+1)m) = (i-k) m.
\]
Karena $p$ membagi $1 + (i+1)m$, bilangan itu tidak mungkin juga membagi $m$
(jika demikian, pembagian pertama akan menyisakan~$1$). Jadi $p$ membagi
$i-k$, karena $p$ membagi $(i-k)m$. Namun, $\left|i-k\right|$ paling besar
bernilai~$n$, dan kita telah memilih $j \geq n$, sehingga ini menyiratkan
bahwa $p \mid m$, sekali lagi suatu kontradiksi. Jadi, tidak ada bilangan
prima yang membagi baik $x_i$ maupun $x_k$. Butir~\olref{less} mudah:
kita mempunyai $y_i < j \leq m < x_i$.

Sekarang mari kita buktikan lema fungsi $\beta$. Ingat bahwa kita dapat
menggunakan $0$, penerus, tambah, kali, $\Char{=}$, proyeksi, dan setiap
fungsi yang didefinisikan darinya dengan menggunakan komposisi dan
pencarian tak berbatas pada fungsi reguler. Kita juga dapat menggunakan
suatu relasi jika fungsi karakteristiknya dapat didefinisikan dengan cara
tersebut. Seperti sebelumnya, dapat kita tunjukkan bahwa relasi-relasi ini
tertutup terhadap kombinasi Boolean dan kuantifikasi terbatas; misalnya:
\begin{align*}
\fn{not}(x) & \defis \Char{=}(x,0)\\
\bmin{x \leq z}{R(x,y)} & \defis \umin{x}{(R(x,y) \lor x = z)}\\
\bexists{x \leq z}{R(x,y)} & \defiff R(\bmin{x \leq z}{R(x,y)}, y)
\end{align*}
Kemudian dapat kita tunjukkan bahwa semua hal berikut juga dapat didefinisikan
tanpa rekursi primitif:
\begin{enumerate}
\item Fungsi pemasangan, $J(x,y) = \frac{1}{2}[(x+y)(x+y+1)] + x$;
% maybe explain more what is going on here, a bit confusing.
\item fungsi-fungsi proyeksi
\begin{align*}
K(z) & = \bmin{x \leq z}{\bexists{y \leq z}{z = J(x,y)}},\\
L(z) & = \bmin{y \leq z}{\bexists{x \leq z}{z = J(x,y)}};
\end{align*}
\item relasi kurang-dari $x < y$;
\item relasi keterbagian $x \mid y$;
% \item $x \tsub y$
% \item $\fn{Prime}(x)$
% \item Assuming $p$ is prime, the relation ``$x$ is a power of $p$'':
% \[
% \bforall{y \leq x}{(y \mid x \lif y = 1 \lor y = x)}.
% \]
\item fungsi $\fn{rem}(x,y)$ yang mengembalikan sisa ketika $y$ dibagi
  dengan~$x$.
\end{enumerate}
Sekarang definisikan
\begin{align*}
\beta^*(d_0,d_1,i) & = \fn{rem}(1+(i+1) d_1,d_0) \text{ dan}\\
\beta(d,i) & = \beta^*(K(d),L(d),i).
\end{align*}
Inilah fungsi yang kita inginkan. Diberikan $a_0,\dots,a_n$ seperti di atas,
misalkan
\[
j = \max(n,a_0+1,\dots,a_n+1),
\]
dan misalkan $d_1 = \lcm(1,\dots,j)$. Menurut \olref{rel-prime} di atas,
kita tahu bahwa $1+d_1$, $1+2 d_1$, \dots, $1+(n+1) d_1$ relatif prima
sepasang-sepasang, dan menurut~\olref{less} bahwa semuanya lebih besar
daripada $a_0,\dots,a_n$. Menurut Teorema Sunzi terdapat nilai~$d_0$
sedemikian sehingga, untuk setiap~$i$,
\[
d_0 \equiv a_i \mod (1+(i+1)d_1)
\]
dan karena itu (karena $d_1$ lebih besar daripada~$a_i$),
\[
a_i = \fn{rem}(1+(i+1)d_1,d_0).
\]
Misalkan $d = J(d_0,d_1)$. Maka, untuk setiap $i \le n$, kita mempunyai
\begin{align*}
\beta(d,i) & =  \beta^*(d_0,d_1,i) \\
& =  \fn{rem}(1+(i+1) d_1,d_0) \\
& =  a_i
\end{align*}
sebagaimana diperlukan. Ini menyelesaikan pembuktian lema fungsi $\beta$.

\begin{prob}
  Tunjukkan bahwa relasi $x < y$, $x \mid y$, dan fungsi
  $\fn{rem}(x,y)$ dapat didefinisikan tanpa rekursi primitif. Anda boleh
  menggunakan $0$, penerus, tambah, kali, $\Char{=}$, proyeksi, serta
  minimisasi berbatas dan kuantifikasi terbatas.
\end{prob}

\end{document}
